\documentclass[]{../../../../tex/classes/styledArticle}
\usepackage{../../../../tex/packages/hebrewSupport}
\usepackage{../../../../tex/packages/mathShortcuts}
\usepackage{../../../../tex/packages/theoremsSupport}
\usepackage{../../../../tex/packages/enfancysections}


%! ~~~ Document ~~~


\author{שחר פרץ}
\title{\textit{ליניארית 2א}}
\date{19 במרץ 2025}
\begin{document}
	\maketitle
	\section{\en{Rules}}
	\begin{multicols}{2}
		\begin{itemize}
			\item תהיו נחמדים, אל תפריעו, תהיו כאן בשביל ללמוד.
			\item לא באמת צריך לדבר בהצבעה אלא אם תהיה בעיה. 
			\item לא נלמד מעבר להיקף החומר הנלמד. 
			\item תרגילי הבית – שבועיים. ליום שני, עולה בשלישי. 
			\item הסיכום שבן משתמש בו יעלה למודל. 
			\item הקורס יותר חישובי. 
			\item הקורס עמוס, תצטרכו לעקוב. לא לפספס משהו ואז  להשלים לקראת תקופת המבחנים. 
			\item האורך המצופה של התרגילים – "ארוך". 
		\end{itemize}
	\end{multicols}
	עומרי המתרגל. יכול להיות שהוא יתקדם קצת יותר מבן. 
	\section{\en{Intoduction}}
	כנגד שלושה נושאים דיברה התורה –
	\begin{enumerate}
		\item \textbf{אופרטורים ליניארים} שיובילו אותנו לצורת ג'ורדן
		\item \textbf{תבניות בי־ליניאריות}
		\item \textbf{מרחבי מכפלה פנימית}
	\end{enumerate}
	בעזרת אלגברה ליניארית אפשר לתאר כל מני דברים גיאומטריים במרחב. עד כה, לא פתחנו את נושא הזוויות. מרחבי מכפלה פנימית הוא הפרמול האלגברי של הגיאומטריה, רק יותר גנרי. 
	
	זה לא בדיוק שלושה שלשים. תבניות בי־ליניאריות נושא קטן, אופרטורים ליניארים נושא גדול ומרחבי מכפלה פנימית פחות גדול. 
	\section{\en{Inner Product}}
	נאמר שישנה פעולה כשהי שנרצה להפעיל. נרצה לקרות מה יקרה לפעולה לאחר שנפעיל אותה מספר פעמים. כפל מטריצות – משהו בגודל של $n^3$. אך, ישנן מטריצות שמאוד קל להעלות בריבוע. 
	\subsection{מטריצות אלכסוניות}
	נאמר ש־$A$ מטריצה מסדר: 
	\[ A = \pms{\lg_1 & & &\\ &\lg_2 && \\ &&\ddots & \\ &&&\lg_n} =: \diag(\lg_1, \lg_2 \dots \lg_n) \]
	ונדבר על ההעתקה: 
	\[ T_A^m(A) = A^m(v) = \pms{\lg_1^m && \\ & \ddots & \\ &&\lg_n^m} \]
	נזכר בסדרת פיבונ'צי. נתבונן בהעתקה הבאה: 
	\[ \pms{a_{n + 2} \\ a_{n + 1}} = \pms{1 & 1 \\ 1 & 0}\pms{a_{n + 1} \\ a_n} \implies \pms{a_n \\ a_{n - 1}} = \pms{1 & 1 \\ 1 & 0}^{n - 1}\pms{1 \\ 0} \]
	(בהנחת איברי בסיס $a_0 = 0, a_1 = 1$). 
	
	מה חבל שלא כיף להכפיל את המטריצה הזו בעצמה המון פעמים. מה נוכל לעשות? ננסה למצוא בסיס שבו $\binom{1 \, 1}{1\, 0}_B = (v_1, v_2)$ ו־$\binom{1 \, 1}{1 \, 0} = P\op \Lambda P$. [המשמעות של $\Lambda$ היא מטריצה לכסינה כלשהי] אז נקבל: 
	\[ \pms{1 & 1 \\ 1 & 0}^{n} = \cl{P\op \Lambda P} = P\op\Lambda^nP \]
	(באינדוקציה – די קל). במקרה כזה יהיה נורא נחמד כי אין בעיה להעלות לכסינה בחזקה. 
	
	הדבר הנחמד הבא שנוכל ליצור הוא צורה ג'ורדנית – מטריצת בלוקים אלכסונית, ככה שנעלה בחזקה את הבלוקים במקום את כל המטריצה. 
	
	\textbf{הגדרה. }\textit{אופרטור ליניארי} (א"ל) הוא ה"ל/טל ממרחב וקטורי $V$ לעצמו. 
	
	מה המשמעות של מטריצה אלכסונית? 
	\[ \Lambda = \diag(\lg_1, \lg_2 \dots \lg_n) \implies \Lambda \pms{\lg_1 \\ 0 \\ \vdots} = \lg_1 \pms{\lg_1 \\ 0 \\ \vdots} \]
	באופן כללי: 
	\[ \Lambda \pms{\vdots \\ 0 \\ \lg_k \\ 0 \\ \vdots} = \lg_k \pms{\vdots \\ 0 \\ \lg_k \\ 0 \\ \vdots} \]
	מה שמוביל אותנו למוטיבציה להגדרה הבאה: 
	
	\textbf{הגדרה. }יהי $T \co V \to V$ א"ל. אז $0 \neq v \in V$ נקרא \textit{וקטור עצמי} של $T$ (ו"ע) אם קיים $\lg \in \F$ כך ש־$T v =\lg v$
	
	\textbf{הגדרה. }$\lg$ מההגדרה הקודמת נקרא \textit{ערך עצמי} (ע"ע) של $T$, המתאים לו"ע $v$. 
	
	\textbf{שאלה. }יהי $T \co \F^n \to \F^n$. נניח ש־$V = (v_1 \dots v_n)$ בסיס של ו"ע של $\F^n$ [תיאורטית יכול להתקיים באופן ריק כי עדיין לא הראינו שקיים בסיס כזה] אז קיימת $P \in M_n(\F)$ הפיכה כך ש־$A$ המקיימת $Tv = Av$ לפי הבסיס הסטנדרטי, אז $[T]^B_B = \diag(\lg_1 \dots \lg_n)$, כאשר $\lg_1 \dots \lg_n$ ע"ע המתאימים לו"ע $v_1 \dots v_n$. 
	
	[למה זו שאלה בכלל?]
	
	"ראיתם את המרחב הומו?"
	כדאי לדעת כי $\Hom(\F^n, \F^n) \approxeq M_{m \times n}(\F)$. מה המשמעות של איזומורפי ($\approxeq$)? בהינתן $A, B$ מבנים אלגברים כלשהם, נסמן $A \approxeq B$ אם קיימת $\phi \co A \to B$ אם קיימת העתקה חח"ע ועל שמכבדת את המבנה. 
	
	\textbf{דוגמה. }אם $V, U$ מ"ו מעל $\F$, הם נקראים איזומורפים אם קיימת $\phi \co V \to U$ חח"ע ועל המקיימת $\forall \lg_1, \lg_2 \in \F \forall v_1, v_2 \in V \co \phi(\lg_1v_1 + \lg_2v_2) = \lg_1\phi(v_1) + \lg_2\phi(v_2)$
	
	כלומר "המרנו" שני מבנים אלגבריים, אבל לא באמת עשינו שום דבר – כל דבר עדיין שומר על התכונות שלו. 
	
	"הספקנו קצת פחות ממה שרציתי אבל כן הספקנו לדבר קצת על דברים פילוסופיים, שהם חשובים". 
	
	
\end{document}
